\documentclass[a4paper,12pt]{article}

\usepackage{fontspec}
\setmainfont{PT Serif}
\usepackage{babel}
\babelprovide[import,main]{russian}
\usepackage{microtype}
\usepackage[margin=2.5cm]{geometry}

\setlength{\parindent}{1.25cm}
\setlength{\parskip}{0pt}

\begin{document}

\begin{center}
    {\Large\bfseries Линейная алгебра}

    \vspace{0.3em}

    {\large Лекция 6}

    \vspace{0.3em}

    \textbf{Линейные отображения}

    \vspace{0.3em}

    {\small 12 сентября 2026 года}
\end{center}

\vspace{1em}

\section{Линейные отображения}

Вводный теоретический материал раздела.

\subsection{Ядро отображения}

Материал первого подраздела.

% Локальные определения, формулы и схемы выбираются из bricks.

\subsection{Образ отображения}

Материал второго подраздела.

\section{Связь с матрицами}

Следующий крупный смысловой блок лекции.

\end{document}
